%%%%%%%%%%%%%%%%%%%%%%%%%%%%%%%%%%%%%%%%%%%%%%%%%%%%%
%			HLAVIČKA								%
%%%%%%%%%%%%%%%%%%%%%%%%%%%%%%%%%%%%%%%%%%%%%%%%%%%%%
\documentclass[openany,12pt]{memoir}

\usepackage[utf8]{inputenc} 
\usepackage[czech]{babel}
\usepackage[T1]{fontenc}
\usepackage[top=1.5cm, bottom=2cm, left=2cm, right=2cm]{geometry}  % --> NASTAVENÍ OKRAJŮ
\usepackage{fancyhdr}
\usepackage{graphicx}
\usepackage{lmodern}
\usepackage{xwatermark}
\usepackage{xcolor}
\usepackage{changepage}




%%%%%% Package na zpěvník
\usepackage[full]{leadsheets}%http://mirrors.nic.cz/tex-archive/macros/latex/contrib/leadsheets/leadsheets_en.pdf   --> dokumentace	
\definesongtitletemplate{empty}{} 
\setchords{
format = \bfseries,   %tučné akordy
minor = {mi},% 
input-notation = {german},%
output-notation = {german}%
}
\definesongtitletemplate{empty}{} 

\newlength{\drop}
\newwatermark[allpages,color=red!50,angle=0,scale=2, xpos=0,ypos=0]{\includegraphics[width=5cm]{obr/pozadi.jpg}} %--> dvojka na pozadí


%%%% Vlastní příkazy
\newcounter{Slokočet}   %Automatické číslování slok
\newcommand{\mezera}{\vspace*{0.5cm}}   %Horizontální odsazení slok
\newcommand{\stred}{5.2cm}   %%% Na zarovnání slok doprostřed, pozn. automatičtější zarovnávání na střed nejde
\newcommand{\refren}{\mezera \noindent \textbf{R:} } %refrén
\newcommand{\sloka}{\mezera \noindent \addtocounter{Slokočet}{1} \arabic{Slokočet}. } 	%sloka, která se automaticky čísluje
\newcommand{\ssloka}{\mezera \noindent}  % vlastní číslo sloky


%%%%%%%%%%%%%%%%%%%%%%%%%%%%%%%%%%%%%%%%%%%%%%%%%%%%%%
%			NÁVOD									 %
%%%%%%%%%%%%%%%%%%%%%%%%%%%%%%%%%%%%%%%%%%%%%%%%%%%%%%
%1. Věci v hlavičce IGNOROVAT
%2. Píseň psát do prostoru mezi \begin{song} a \end{song}
%3. další řádek se značí dvěma odsazeníma (= dvakrát stisknout enter)
%4. \refren vždy na začátku refrenu a \sloka na začátek sloky (automaticky se čísluje)
%  \ = alt gr + q ; [/] = alt gr f/g ; {/} = alt gr + b/n; ^ = alt gr + 3 + mezera
%Cokoliv napíšete do ^{  } se bude brát jako akord
%když se toto bude dotýkat nějakého slova (nebude mezi tím a slovem mezera)
%tak se akordy zjeví nad slovem, ale pište to před slovo
%Když se to nedotýká slova, tak akord lítá ve vzduchu a vytiskne se větší mezera
%První možnost je asi preferovanější
%5. Akordy stačí psát jen do první sloky, když se nezmění -- kytaristi to zvládnou



%<++>
%\usepackage{subfiles}
%</++>

\begin{document}

\pagestyle{empty}
%\addtocontents{toc}{\protect\thispagestyle{empty}}
%\tableofcontents \thispagestyle{empty}\newpage
\newgeometry{top=1.5cm, bottom = 0cm, left = 2cm, right = 2cm}


\documentclass[../main.tex]{subfiles}

\chapter{God's gonna cut you down}

\noindent\hspace{0.15\linewidth}\begin{minipage}{0.7\linewidth}
\begin{song}{title=Johnny Cash - Gods gonna cut you down}
\begin{verse}[numbered]
You^{Am} can run on for a long time, \\
Run^{Am} on for a long time, \\
Run^{Am} on for a long time, \\
Sooner, or later^{Dm}, God'll cut you down^{Am}. \\
Sooner, or later^{Dm}, God'll cut you down^{Am}. \\
\end{verse}
\begin{verse}[numbered]
Go and tell^{Am} that long tongue liar, \\
Go and tell^{Am} that midnight rider, \\
Tell the rambler^{Am}, the gambler, the back biter, \\
Tell^{C} 'em that God's^{Em} gonna cut^{Am} 'em down. \\
Tell^{C} 'em that God's^{Em} gonna cut^{Am} 'em down. \\
\end{verse}
\begin{verse}[numbered]
Well my^{Am} goodness gracious, let me tell you the news. \\
My heads^{Am} been wet with the midnight dew. \\
I've^{Am} been down on bended knee, \\
talkin^{Am} to the man from Galiee. \\
He spoke^{Am} to me in a voice so sweet, \\
I thought^{Am} I heard the shuffle of angels feet.  \\
He called^{Am} my name and my heart stood still,  \\
When He^{Am} said^{Mute} \uv{John go do my will}  \\
\end{verse}
\begin{verse}[numbered]
Go and tell^{Am} that long tongue liar, \\
Go and tell^{Am} that midnight rider, \\
Tell the rambiler^{Am}, the gambler, the back biter, \\
Tell^{C} 'em that God's^{Em} gonna cut 'em down^{Am}. \\
Tell^{C} 'em that God's^{Em} gonna cut 'em down^{Am}. \\
\end{verse}
\end{song}
\end{minipage}

\newpage
\vspace*{4.5cm}
\noindent\hspace{0.15\linewidth}\begin{minipage}{0.7\linewidth}
\begin{song}{title=Johnny Cash - Gods gonna cut you down}
\begin{verse}
\item[5.]You can run^{Am} on for a long time, \\
Run^{Am} on for a long time, \\
Run^{Am} on for a long time, \\
Sooner^{C}, or later,^{Em} God'll cut^{Am} you down. \\
Sooner^{C}, or later,^{Em} God'll cut^{Am} you down. \\
You may throw^{Am} your rock, hide your hand, \\
Workin^{Am} in the dark against your fellow man. \\
But as^{Am} sure as God made black and white, \\
What's done^{Am} in the dark, will be brought to the light. \\
\end{verse}
\begin{verse}
\item[6.]You can run^{Am} on for a long time, \\
Run^{Am} on for a long time, \\
Run^{Am} on for a long time, \\
Sooner^{C}, or later,^{Em} God'll cut^{Am} you down. \\
Sooner^{C}, or later,^{Em} God'll cut^{Am} you down. \\
\end{verse}
\begin{verse}
\item[7.]Go and tell^{Am} that long tongue liar, \\
Go and tell^{Am} that midnight rider, \\
Tell the rambler^{Am}, the gambler, the back biter, \\
Tell^{C} 'em that God's^{Em} gonna cut^{Am} 'em down. \\
Tell^{C} 'em that God's^{Em} gonna cut^{Am} 'em down^{Am}. \\
Tell^{C} 'em that God's^{Em} gonna cut^{Am} 'em down^{Am}. 
\end{verse}
\end{song}
\end{minipage}

\newpage %!!! Akordy dát nad nad slova, vyřadit opakující se akordy
\end{document}
